%\documentclass[12pt,oneside]{amsart}
\documentclass{scrartcl}

%\documentclass{article}

\usepackage{amsthm,amsmath,amssymb}
\usepackage[numbers]{natbib}

%\usepackage[colorlinks]{hyperref}
\usepackage{hyperref}
\hypersetup{
    colorlinks,%
    citecolor=black,%
    filecolor=black,%
    linkcolor=black,%
    urlcolor=black
}
\usepackage{hypernat}
\usepackage[top=2.5cm, bottom=2.5cm, left=2.8cm,
right=2.8cm]{geometry} 

%\setlength{\parskip}{0pt}
%\setlength{\parsep}{0pt}
%\setlength{\headsep}{0pt}
%\setlength{\topskip}{0pt}
%\setlength{\topmargin}{0pt}
%\setlength{\topsep}{0pt}
%\setlength{\partopsep}{0pt}

\linespread{1}

%\usepackage{marvosym}

%\textwidth = 6.5 in \textheight = 9.2 in \oddsidemargin = 0.0 in
%\evensidemargin = 0.0 in \topmargin = -0.0 in \headheight = 0.0 in
%\headsep = 0.0in \parskip = 0.0in \parindent = 0.0in
%\def\baselinestretch{0.871}


\newtheorem{theorem}{Theorem}[section]
\newtheorem{proposition}[theorem]{Proposition}
\newtheorem{problem}[theorem]{Problem}
\newtheorem{fact}[theorem]{Fact}
\newtheorem{lemma}[theorem]{Lemma}
\newtheorem{defn}[theorem]{Definition}
\newtheorem{corollary}[theorem]{Corollary}
\newtheorem{remark}{Remark}[section]
\newtheorem{example}[theorem]{Example}

\newcommand{\E}{{\mathbb E}}
\newcommand{\se}{{\mathcal{E}}}
\newcommand{\R}{{\mathbb R}}
\renewcommand{\P}{{\mathbb P}}
\newcommand{\ac}{{\mathcal{A}}}
\newcommand{\A}{{\mathcal{A}}}
\newcommand{\B}{{\mathcal{B}}}
\newcommand{\C}{{\mathcal{C}}}
\newcommand{\M}{{\mathcal{M}}}
\newcommand{\Mf}{{\mathfrak{M}}}
\newcommand{\T}{{\mathcal{T}}}
\newcommand{\Z}{{\mathcal{Z}}}
\newcommand{\K}{{\mathcal{K}}}
\newcommand{\lc}{{\mathcal{L}}}
\newcommand{\Ps}{{\mathcal{P}}}
\newcommand{\G}{{\mathcal{G}}}
\newcommand{\pa}{{\mathring{p}}}
\newcommand{\F}{{\cal F}}
\newcommand{\samp}{{\mathcal{X}}}
\newcommand{\X}{{\mathcal{X}}}
\newcommand{\hi}{{\hat{\Phi}}}
\newcommand{\data}{{\text{data}}}

\newcounter{rcnt}[section]
\renewcommand{\thercnt}{(\roman{rcnt})}

\def\qt#1{\qquad\text{#1}}

\def\argmin{\mathop{\rm argmin}}
\def\argmax{\mathop{\rm argmax}}


\setlength{\parskip}{1.4 \medskipamount}
%\sloppy
%\linespread{1.3}

\begin{document}

\title{STAT 153 \& 248 - Time Series  \\
  Lecture Two} 
\subtitle{Fall 2026, UC Berkeley}
\author{Aditya Guntuboyina}

\date{01 September 2026}

\maketitle


We shall start the first topic of the course: Linear Regression. We
start our discussion with simple linear regression (where there is a
single covariate), and then extend to multiple linear regression
(where there are multiple covariates). 

\section{Simple Linear Regression}
We want to learn the relationship between two variables $y$ and
$x$, with the aim of predicting $y$ given the value of $x$. $y$ is
called the response variable, and $x$ is called the covariate. For
example (this was one of the original applications of regression), $y$
denotes the height of an adult man, and $x$ denotes the height of
their father. The linear regression model assumes that $y$ is related
to $x$ via the equation:
\begin{equation}\label{lin_eq}
  y = \beta_0 + \beta_1 x + \epsilon
\end{equation}
where $\beta_0$ and $\beta_1$ are parameters, and $\epsilon$ denotes
an error term which captures deviations of $y$ from the assumed
equation $\beta_0 + \beta_1 x$. The parameters $\beta_0$ and $\beta_1$
can be interpreted as follows: $\beta_0$ denotes the value of $y$ when
$x = 0$ and $\beta_1$ represents the change in $y$ when $x$ changes by
one unit. 

We observe data $(x_1, y_1), \dots, (x_n, y_n)$ on the covariate and
response variables corresponding to $n$ instances (in the height
example, we have data on heights for $n$ father-son pairs). Writing
the equation \eqref{lin_eq} for each individual pair $(x_i, y_i)$ we
get
\begin{equation}\label{linmod}
  y_i = \beta_0 + \beta_1 x_i + \epsilon_i. 
\end{equation}
The observed data $(x_1, y_1), \dots, (x_n, y_n)$ will be used to
obtain estimates $\hat{\beta}_0$ and $\hat{\beta}_1$ (as well as
uncertaintly quantification) for the parameters $\beta_0$ and
$\beta_1$. After obtaining these estimates, one can predict the value
of the response variable for a possibly new covariate value
$x_{\text{new}}$ by $\hat{\beta}_0 + \hat{\beta}_1
x_{\text{new}}$.

We will implement the equations for obtaining $\hat{\beta}_0$ and
$\hat{\beta}_1$ from the 
observed data $(x_1, y_1), \dots, (x_n, y_n)$ using the Python library
\texttt{statsmodels}. We will also study the math behind
this process. 

To apply linear regression, we need data on both $y$ and $x$. In the
time series context, the observed data is $y_1, \dots, y_n$ which
represent observations for a single variable $y$. There is no
additional data on another variable $x$. In order to apply regression
methods to time series, we need to create a covariate variable
$x$. There are two main ways of doing it:
\begin{enumerate}
\item \textbf{Time as covariate}: Here we take the time index as the
  covariate $x$.  For example, in the time series dataset on the
population of the United States for each month from January 1959 to
December 2024: $n$ denotes the total number of data points, $x_i = i$
and $y_i$ denotes the observed population data for the $i^{\text{th}}$
month (first month is January 1959, second month is February 1959 and
so on).
\item \textbf{Lagged $y$ as covariate}: Here we take $x_i =
  y_{i-1}$. In other words, the covariate equals the response at the
  previous time point. This kind of regression is called Lagged
  Regression or, more commonly, AutoRegression. 
\end{enumerate}
For now, we shall focus on the first kind of regression (time as
covariate). We shall study AutoRegression in more detail later.

\section{Estimation of $\beta_0$ and $\beta_1$}

\subsection{Least Squares Estimates}

The estimates of $\beta_0$ and $\beta_1$ reported by standard
libraries (such as \texttt{statsmodels}) are obtained using the method
of least squares. This involves minimizing the sum of squares
criterion:
\begin{equation}\label{ss}
  S(\beta_0, \beta_1) = \sum_{i=1}^n \left(y_i - \beta_0 - \beta_1 x_i
  \right)^2 
\end{equation}
over all values of $\beta_0$ and $\beta_1$. It is left as an exercise
to verify that: 
\begin{equation*}
  \hat{\beta}_0 = \bar{y} - \hat{\beta}_1 \bar{x}  ~~ \text{ and } ~~
  \hat{\beta}_1 = \frac{\sum_{i=1}^n (y_i - \bar{y})(x_i -
    \bar{x})}{\sum_{i=1}^n (x_i - \bar{x})^2}, 
\end{equation*}
where
\begin{equation*}
  \bar{y} = \frac{y_1 + \dots + y_n}{n} ~~ \text{ and } ~~   \bar{x} =
  \frac{x_1 + \dots + x_n}{n}. 
\end{equation*}

\subsection{MLE under Normality of Errors}
Suppose we assume that the error terms $\epsilon_1, \dots, \epsilon_n$ in
\eqref{linmod} are i.i.d normal with mean zero and some variance
$\sigma^2$:
\begin{equation}\label{errgauss}
 \epsilon_1, \dots, \epsilon_n \overset{\text{i.i.d}}{\sim} N(0,
 \sigma^2). 
\end{equation}
Then the least squares estimates of $\beta_0$ and $\beta_1$ coincide
with the Maximum Likelihood Estimates (MLEs). 

Another way of writing the model \eqref{linmod} and \eqref{errgauss}
is: 
\begin{equation*}
  y_i \overset{\text{independent}}{\sim} N(\beta_0 + \beta_1 x_i, 
  \sigma^2). 
\end{equation*}
To obtain the MLEs of the parameters ($\beta_0, \beta_1$ as well as
$\sigma$), we need to write the likelihood function and then maximize
it. The likelihood function is the joint density of the data for fixed
values of the parameters $\beta_0, \beta_1, \sigma$: 
\begin{align}
  f_{y_1, \dots, y_n \mid \beta_0, \beta_1, \sigma}(y_1, \dots, y_n) &=
  \prod_{i=1}^n \frac{1}{\sqrt{2 \pi} \sigma} \exp \left(-\frac{(y_i -
                                                                       \beta_0
                                                                       -
                                                                       \beta_1
                                                                       x_i)^2}{2
                                                                       \sigma^2}
                                                                       \right)
  \nonumber \\
  &= (2 \pi)^{-n/2} \sigma^{-n} \exp \left(-\frac{1}{2 \sigma^2}
    \sum_{i=1}^n (y_i - \beta_0 - \beta_1 x_i)^2 \right) \nonumber \\
  &= (2 \pi)^{-n/2} \sigma^{-n} \exp \left(-\frac{S(\beta_0,
    \beta_1)}{2 \sigma^2} \right) \label{like}
\end{align}
where $S(\beta_0, \beta_1)$ is the sum of squares \eqref{ss}. 
\begin{equation*}
  S(\beta_0, \beta_1) :=
    \sum_{i=1}^n (y_i - \beta_0 - \beta_1 x_i)^2. 
  \end{equation*}
To write this likelihood, we are assuming that $x_1, \dots, x_n$ are
fixed. This assumption is fine if $x_i = i$ (regression with time as
covariate) but not strictly true when $x_i = y_{i-1}$
(auto-regression). We shall see how it is still approximately true in
the case of AutoRegression later.

Maximization of the likelihood is a three variable optimization
problem (the variables being $\beta_0, \beta_1, \sigma$). The optimal
values of $\beta_0$ and $\beta_1$ in this problem coincide with the
least squares estimate. We shall see why in the next lecture. 
\end{document}

%%% Local Variables:
%%% mode: latex
%%% TeX-master: t
%%% End:
